% sec2a.tex — c04｜原文 p8–12｜§2 超共形指标（开头含 §1 末尾的习题 1.2 与两段收尾）
% 编号公式 (2.1)–(2.16) 共 16 条；编号环境：习题 1.2、2.1；表 1（复刻为 tabular）。
% 结构说明：分块表把 p8 的习题 1.2 归入本文件；原文中它在 §2 标题之前，故编译后会落在
% §2 标题之下。此处按分块表施工，已在块末批注与回报中列为未决。
% [原文 p.8]
\begin{exercise}[1.2]
在 $x$ 展开中以及用 Cauchy 定理对 $\SO(2)$ 计算这个积分。
\end{exercise}

再次，单字母指标为 $i_B(x,a)=0$ 且 $i_F(x,a)=0$。这意味着 $I(x,a)=PE[i_B(x)-i_F(x)]=1$。注意，在指标中只能开启那些与超荷对易的荷的逸度。这是因为指标的稳健性依赖于非零能态上的玻色—费米抵消。若开启与 $Q$ 不对易的荷的逸度，抵消将不成立，指标在量子修正下也不会受到保护。

到目前为止，我们还没有遇到具有非平庸指标的理论，但我们已经发展出在指标非平庸时有助于计算它的工具。现在我们将把注意力转向超共形指标。

四维超共形场论的指标定义为该理论在径向量子化下的 Witten 指标。取 $Q\equiv Q_-$ 为其一 Poincaré 超荷，$Q^\dagger=S_+$ 为共轭的共形超荷。概略地说，指标定义为 [12–14]
\[
\sI(\mu_i)=\Tr(-1)^F x^{\delta}\mu_i^{\mathcal{M}_i} ,
\tag{2.1}
\]
其中迹取在该理论于 $S^3$ 上的 Hilbert 空间上，$\delta\equiv\frac12\{Q,Q^\dagger\}$。由于 $\delta>0$ 的态以玻色/费米对的形式出现，只有 $\delta=0$ 的态有贡献，且指标与 $x$ 无关。如前所述，费米—玻色抵消只有在额外的逸度 $\mu_i$ 仅对与 $Q$ 对易的那些荷 $\mathcal{M}_i$ 开启时才成立。与量子力学的情形不同，该理论有无穷多个 $\delta=0$ 的态。不仅如此，正如我们很快将看到的，该理论还有无穷多个 $\delta=0$ 的「单字母」。逸度 $\mu_i$ 的引入既是为了正则化这一发散，也是为了实现更精细的计数。

对 $\sN=1$，超荷为 $\{Q_\alpha, S^\alpha\equiv Q^{\dagger\alpha}, \tilde Q_{\dot\alpha}, \tilde S^{\dot\alpha}\equiv \tilde Q^{\dagger\dot\alpha}\}$，其中 $\alpha=\pm$ 与 $\dot\alpha=\dot\pm$ 分别是 $\SU(2)_1$ 与 $\SU(2)_2$ 的指标，而 $\SU(2)_1\times \SU(2)_2=\mathrm{Spin}(4)$ 是 $S^3$ 的等距群。相关的反对易关系为
% [原文 p.9]
\[
\{Q_\alpha, Q^{\dagger\beta}\}=\Delta+2M_\alpha^\beta+\frac32 r ,
\tag{2.2}
\]
\[
\{\tilde Q_{\dot\alpha}, \tilde Q^{\dagger\dot\beta}\}=\Delta+2\tilde M_{\dot\alpha}^{\dot\beta}-\frac32 r ,
\tag{2.3}
\]
其中 $\Delta$ 是共形维数，$M_\alpha^\beta$ 与 $\tilde M_{\dot\alpha}^{\dot\beta}$ 是 $\SU(2)_1$ 与 $\SU(2)_2$ 生成元，$r$ 是 $\U(1)_r$ R 对称性的生成元。在我们的约定中，$Q$ 的 $r=-1$，而 $\tilde Q$ 的 $r=+1$；当然，dagger 操作会翻转 $r$ 的符号。

取超荷 $Q\equiv Q_-$ 来定义指标，我们得到 $\delta=\Delta-2j_1+\frac32 r$。$Q$ 在超共形代数 $\SU(2,2|\sN)$ 中的对易子为 $\SU(2,1|\sN-1)$。对 $\sN=1$，对易子为 $\SU(2,1)$。它的秩为 2。我们可以在这个空间中选择任意一组荷作为 $\mathcal{M}_i$ 来精细化指标。为了与已有文献一致，我们取其为 $\frac13(\Delta+j_1)\pm j_2$。于是指标定义为
\[
\begin{aligned}
\sI(p,q)&\equiv\Tr(-1)^F p^{\frac13(\Delta+j_1)+j_2}q^{\frac13(\Delta+j_1)-j_2}\\
&=\Tr(-1)^F p^{j_1-\frac12 r+j_2}q^{j_1-\frac12 r-j_2},\qquad \delta=\Delta-2j_1+\frac32 r ,
\end{aligned}
\tag{2.4}
\]
在第二个等号中，我们利用只有 $\delta=0$ 的态对指标有贡献这一事实，把荷 $\frac13(\Delta+j_1)$ 换成了 $j_1-\frac12 r$。这样做是因为 $r$ 荷比共形维数更容易确定。

\subsection{手征多重态}

我们从单个自由手征多重态 $\Phi$ 的理论开始。它满足条件 $\bar Q_{\dot\alpha}\Phi=0$。手征多重态的超空间展开为
\[
\Phi=\phi+\sqrt2\,\theta\psi+\ii\,\theta^\dagger\bar\sigma^\mu\theta\,\partial_\mu\phi .
\tag{2.5}
\]
复共轭是一个反手征场，即 $Q_\alpha\bar\Phi=0$。$\phi$ 的 $r$ 荷为 $2/3$（你知道为什么吗？）。

在径向量子化中，理论的 Hilbert 空间同构于局域算符的谱。这些算符由手征多重态中的场及其导数构造。这意味着，量子力学中产生算符的对应物是场及其导数。它们是
\[
\phi,\ \partial_\mu\phi,\ \partial_\mu\partial_\nu\phi,\ \ldots,\qquad
\psi,\ \partial_\mu\psi,\ \partial_\mu\partial_\nu\psi,\ \ldots
\tag{2.6}
\]
及其共轭。这里导数以所有可能的缩并作用，但带有来自运动方程的约束 $\partial_\mu\partial^\mu\phi=0$。类似地，费米子也满足运动方程 $\partial_{\alpha\dot\alpha}\psi^\alpha=0$。有了这些约束，我们在表 1 中列出字母并枚举「单字母指标」。由此，我们可以容易地计算单字母指标，
\[
i_\phi(p,q)=\frac{(pq)^{\frac13}-(pq)^{\frac23}}{(1-p)(1-q)} .
\tag{2.7}
\]
% [原文 p.10]
\noindent\textbf{表 1.} $\sN=1$ 手征多重态的「字母」及其对指标的贡献。其中 $\delta=\Delta-2j_1+\frac32 r$。先验地，我们必须考虑自由运动方程 $\partial\psi=0$ 与 $\Box\phi=0$，它们对可能的词施加约束；但我们看到，在此情形下运动方程的 $\delta\neq0$，因此它们对单字母指标没有贡献。最后有两个时空导数对指标有贡献（即 $\delta=0$），它们在场上形成的塔负责单字母指标的分母。

\begin{center}
\begin{tabular}{C{4.6em} C{1.4em} C{1.4em} C{1.4em} C{1.4em} C{1.4em} C{2.2em}}
\toprule
字母 & $\Delta$ & $j_1$ & $j_2$ & $r$ & $\delta$ & $\sI$ \\
\midrule
$\phi$ & $1$ & $0$ & $0$ & $\frac23$ & $2$ & \\
$\psi_+$ & $\frac32$ & $\frac12$ & $0$ & $-\frac13$ & $0$ & $-(pq)^{\frac23}$ \\
$\psi_-$ & $\frac32$ & $-\frac12$ & $0$ & $-\frac13$ & $2$ & \\
$\partial^\alpha_{\dot+}\psi_\alpha$ & $\frac52$ & $0$ & $\frac12$ & $-\frac13$ & $2$ & \\
$\partial^\alpha_{\dot-}\psi_\alpha$ & $\frac52$ & $0$ & $-\frac12$ & $-\frac13$ & $2$ & \\
$\Box\phi$ & $3$ & $0$ & $0$ & $\frac23$ & $4$ & \\
\midrule
$\bar\phi$ & $1$ & $0$ & $0$ & $-\frac23$ & $0$ & $(pq)^{\frac13}$ \\
$\bar\psi_{\dot+}$ & $\frac32$ & $0$ & $\frac12$ & $\frac13$ & $2$ & \\
$\bar\psi_{\dot-}$ & $\frac32$ & $0$ & $-\frac12$ & $\frac13$ & $2$ & \\
$\partial^{\dot\alpha}_{+}\bar\psi_{\dot\alpha}$ & $\frac52$ & $\frac12$ & $0$ & $\frac13$ & $2$ & \\
$\partial^{\dot\alpha}_{-}\bar\psi_{\dot\alpha}$ & $\frac52$ & $-\frac12$ & $0$ & $\frac13$ & $4$ & \\
$\Box\bar\phi$ & $3$ & $0$ & $0$ & $-\frac23$ & $2$ & \\
\midrule
$\partial_{+\dot+}$ & $1$ & $\frac12$ & $\frac12$ & $0$ & $0$ & $p$ \\
$\partial_{-\dot+}$ & $1$ & $-\frac12$ & $\frac12$ & $0$ & $0$ & $q$ \\
$\partial_{+\dot-}$ & $1$ & $\frac12$ & $-\frac12$ & $0$ & $2$ & \\
$\partial_{-\dot-}$ & $1$ & $-\frac12$ & $-\frac12$ & $0$ & $2$ & \\
\bottomrule
\end{tabular}
\end{center}

分母编码了两个 $\delta=0$ 的时空导数的作用。如果手征多重态在某个整体对称性的表示 $R$ 中变换，那么我们可以用 Cartan 生成元的逸度来精细化单字母指标。注意，如果手征多重态在表示 $R$ 中变换，则反手征多重态在共轭表示 $\bar R$ 中变换。
\[
i_\phi(a_i;p,q)=\frac{(pq)^{\frac13}\chi_{\bar R}(a_i)-(pq)^{\frac23}\chi_R(a_i)}{(1-p)(1-q)} .
\tag{2.8}
\]
\begin{exercise}[2.1]
相对于 $\tilde Q_{\dot+}$ 定义指标，并仿照类似的表格计算 $\Phi$ 的单字母指标。
\end{exercise}
全指标通过齐次指数化得到。对在 $\U(1)$ 对称性下荷为 1 的单个多重态
% [原文 p.11]
，其指标为
\[
\sI_\phi(a;p,q)=\Gamma\big((pq)^{\frac13}a^{-1};p,q\big),\qquad
\Gamma(z;p,q)\equiv\prod_{i,j=1}^{\infty}\frac{1-p^{i+1}q^{j+1}/z}{1-p^iq^jz} .
\tag{2.9}
\]
特殊函数 $\Gamma(z;p,q)$ 称为椭圆 gamma 函数 [15]，在数学文献中作为通常 Gamma 函数的「椭圆」推广而被研究。椭圆 gamma 函数在超共形指标理论中起重要作用，因为手征多重态是四维任何超共形拉氏量不可或缺的组成部分。我们几乎总是略去椭圆 Gamma 函数的自变量 $(p,q)$，并记 $\Gamma(z;p,q)\to\Gamma(z)$。如果手征多重态在表示 $R$ 中变换，则相应的指标为
\[
\sI_\phi(a_i;p,q)=\prod_{\rho\in R}\Gamma\big((pq)^{\frac13}a^\rho\big) .
\tag{2.10}
\]
如果手征多重态参与超势相互作用，那么它的 R 荷可以不同于 $2/3$。它由要求超势具有 R 荷 2 来确定。如果理论具有阿贝尔整体对称性，这一约束不能完全固定 R 荷，因为这些对称性可以与这个假定的 R 荷混合。在那种情形下，R 荷由要求 R 对称性与任何阿贝尔整体对称性都没有混合反常来确定。这一程序称为「a-极大化」[16,17]。荷为 $r$ 的手征多重态的单字母指标，是遍及自由情形下对指标有贡献的字母求和得到的；但这些字母按其 R 荷对指标有贡献。这称为 Romelsberger 处方 [13,14]。这一奇特处方后来通过在 $S^3\times S^1$ 上沿 RG 流保持超对称而被理解 [18]。

由表 1，手征多重态中对指标有贡献的字母是 $\psi_+$ 与 $\bar\phi$。若手征多重态的 R 荷为 $r$，则这两个字母的 R 荷分别为 $r-1$ 与 $-r$。相应地，它们将分别贡献 $-(pq)^{1-\frac12 r}$ 与 $(pq)^{\frac12 r}$。这给出 R 荷为 $r$ 的手征多重态的指标为
\[
i_{\phi_r}(p,q)=\frac{(pq)^{\frac12 r}-(pq)^{1-\frac12 r}}{(1-p)(1-q)} .
\tag{2.11}
\]
多粒子指标为
\[
\sI_{\phi_r}(a;p,q)=\Gamma\big((pq)^{\frac12 r}a\big) .
\tag{2.12}
\]

\paragraph{例子} 有了这些工具，我们已经可以计算一些只涉及手征多重态的物理理论的指标。本小节的例子取自综述 [8]。

\paragraph{有质量理论：} 考虑带质量项的单个手征多重态理论，即 $W(\Phi)=m\Phi^2$。这把手征场的 R 荷固定为 1。
\[
\sI_{\phi_1}(p,q)=\Gamma\big((pq)^{\frac12}\big)=1 .
\tag{2.13}
\]
% [原文 p.12]
这与物理预期一致：有质量理论应当只有一个超对称基态。

作为轻微的推广，可以考虑两个手征多重态的理论，其质量型超势为 $W(\Phi_i)=m\Phi_1\Phi_2$。这个超势不足以完全固定 $\Phi_{1,2}$ 的 R 荷，但它确实告诉我们 $r_1+r_2=2$。仅凭这一约束
\[
\sI(a;p,q)=\Gamma\big((pq)^{\frac12 r}a\big)\Gamma\big((pq)^{1-\frac12 r}a^{-1}\big)=1 .
\tag{2.14}
\]
这也在预料之中，因为该超势使两个手征多重态都获得能隙，从分量形式写出的拉氏量可以清楚地看出这一点。

\paragraph{F 项超对称破缺} 考虑一个简单的 O'Raifeartaigh 型模型，其超势为线性 $W(\Phi)=\eta\Phi$。该理论自发破缺超对称，因为该超势给基态以非零能量。让我们看看这在指标中如何体现。该超势把 $\Phi$ 的 R 荷固定为 2，且不允许任何整体对称性。
\[
\sI_{\phi_2}(p,q)=\Gamma(pq)=0 .
\tag{2.15}
\]
指标中的零来自 $\psi_+$ 的贡献变为 $-(pq)^{1-\frac12 r}\to-1$。这意味着带 $\psi_+$ 的态 $\psi_+|0\rangle$ 与 $|0\rangle$ 具有相同的能量。换言之，基态以玻色—费米对的形式出现，因此不受保护。这也意味着 $\psi_+$ 应被解释为超对称自发破缺的 Goldstino 模。任何自发破缺超对称的 O'Raifeartaigh 型模型总有一个荷为 2 的中性手征多重态。这导致 $\sI=0$。

\paragraph{非平庸手征环} 具有非平庸手征环的最简单模型是单个手征多重态的理论，其超势为 $W(\Phi)=\Phi^{n+1}$。手征环关系为 $W'(\Phi)=0=\Phi^n$。超势意味着场 $\Phi$ 具有 R 荷 $2/(n+1)$，并且它关于一个离散对称性 $\mathbb{Z}_{n+1}$ 带荷 1。我们可以开启味对称性的逸度 $g$ 来计算指标。由于味对称性只是离散的 $\mathbb{Z}_{n+1}$，逸度满足 $g^{n+1}=1$。单字母指标为
\[
i(g;p,q)=\frac{(pq)^{\frac{1}{n+1}}g-(pq)^{1-\frac{1}{n+1}}g^{-1}}{(1-p)(1-q)}
=\frac{(pq)^{\frac{1}{n+1}}g-\big((pq)^{\frac{1}{n+1}}g\big)^n}{(1-p)(1-q)} .
\tag{2.16}
\]
来自费米字母 $\psi_+$ 的贡献恰好是玻色字母贡献的 $n$ 倍，且符号相反。这与手征环关系 $\Phi^n=0$ 相符。

% [Tutor Note: 术语：本节的「指标」$\sI$ 即四维 SCFT 的 Witten 指标在径向量子化下的版本；与 §1 的配分函数 $Z(x)$ 不同，它能对超共形代数做精细计数。]
% [Tutor Note: 动机：$\delta\equiv\frac12\{Q,Q^\dagger\}$ 扮演「能量」；$\delta>0$ 的态成对抵消，故指标只数 BPS 态，且与 $x$ 无关。]
% [Tutor Note: 易错：超共形指标中 $x^\delta$ 的 $\delta$ 已换成 $\Delta-2j_1+\frac32 r$，不要与 §1 的哈密顿量 $H$ 混同。]
% [Tutor Note: 易错：$\mu_i$ 只能对与 $Q$ 对易的荷 $\mathcal{M}_i$ 开启，否则玻色—费米抵消失效，指标不再受保护。]
% [Tutor Note: 术语：$S^\alpha\equiv Q^{\dagger\alpha}$ 是共形超荷；$\alpha$ 与 $\dot\alpha$ 分别属于 $\SU(2)_1$、$\SU(2)_2$，合成 $\mathrm{Spin}(4)$，不要与 R 对称性混淆。]
% [Tutor Note: 动机：$Q$ 的对易子 $\SU(2,1|\sN-1)$ 的秩决定可开启的逸度个数；$\sN=1$ 时秩 2，故 $\sI$ 依赖 $p,q$。]
% [Tutor Note: 易错：(2.4) 第二个等号只在 $\delta=0$ 的态上成立，它把较难确定的共形维数 $\Delta$ 换成 $r$ 荷。]
% [Tutor Note: 术语：$\Gamma(z;p,q)$ 是椭圆 gamma 函数，不要与普通 gamma 函数 $\Gamma(z)$ 混淆；后者是 $p,q\to0$ 的极限。]
% [Tutor Note: 动机：表 1 的分母 $(1-p)(1-q)$ 来自两个 $\delta=0$ 的时空导数 $\partial_{+\dot+}$、$\partial_{+\dot-}$ 的塔。]
% [Tutor Note: 易错：表 1 第三组第 2、3 行的 $\delta$ 与 $\sI$ 列疑有笔误：按 $\delta=\Delta-2j_1+\frac32 r$，$\partial_{-\dot+}$ 的 $\delta$ 应为 $2$、$\partial_{+\dot-}$ 的 $\delta$ 应为 $0$（$q$ 应随 $\delta=0$ 行）；译文照录，已登记 ERRATA。]
% [Tutor Note: 易错：表 1 只列「单字母」；词的贡献要再把单字母指标做齐次指数化（(2.9) 的 $\Gamma$ 就是这一塔的求和）。]
% [Tutor Note: 易错：自由手征多重态 $r=2/3$；有超势时 $r$ 由 $W$ 的 R 荷 2 与 a-极大化决定，不能一律取 $2/3$。]
% [Tutor Note: 术语：Romelsberger 处方＝用自由理论的字母、但按相互作用理论的 R 荷赋权；文献 [13,14] 提出，[18] 用 $S^3\times S^1$ 上的超对称流解释。]
% [Tutor Note: 易错：$\sI=0$ 不是「没有态」，而是基态成玻色—费米对、超对称自发破缺；此时 $\psi_+$ 是 Goldstino。]
% [Tutor Note: 易错：$\mathbb{Z}_{n+1}$ 是离散味对称性，逸度满足 $g^{n+1}=1$，不能当作连续 $\U(1)$ 的逸度使用。]
% [Tutor Note: 文献：[15] 椭圆 gamma 函数；[16,17] a-极大化；[8] 综述中的例子。]
% [Tutor Note: 结构提醒：习题 1.2 原文在 §1 末、§2 标题之前，但按分块表归入本文件，编译后会出现在 §2 标题之下。]
% [Tutor Note: 语病：原文 p.11 「the index it studied」缺系动词，应为「the index it is studied」或「the index studied」；照译并登记 ERRATA。]
% [Tutor Note: 语病：原文 p.12 「the this superpotential」冠词重复，应为「this superpotential」；已登记 ERRATA。]
